\section{Introduction}
\label{sec:intro}

\IEEEPARstart{R}{ecycling}, composting and safe disposal of waste are possible only when the different materials are kept apart. In India, the Solid Waste Management Rules, 2016 make every waste generator responsible for segregating waste at source into bio-degradable, non-bio-degradable and domestic hazardous streams (Rule~4), and prescribe green, white and black bins for bio-degradable, recyclable and other waste (Rule~15(h))~\cite{moefcc2016solid}. A camera-based system that can identify the material of each item would help this segregation at a smart bin, on a sorting line or at a collection point. The use of artificial intelligence in solid waste management is by now well established~\cite{abdallah2020artificial,fang2023artificial}.

Most of the vision work on waste sorting, however, deals with a simpler problem than the one described above. A typical study takes a photograph of one object, usually on a plain background, and assigns it a class, as in TrashNet~\cite{yang2016classification} and the classifiers built on it~\cite{kaya2023classification}. A heap of waste, a conveyor belt or a street contains many items at once. Some of them touch or hide each other, and every item has to be located as well as named, which makes it an object-detection problem. Majchrowska et al.~\cite{majchrowska2022deep}, whose work was the starting point for ours, first localize litter with one network and then classify each detected crop with a second one. Multi-object waste detection therefore exists, but it remains a smaller body of work than single-object classification, and several detection studies still depend on single-object studio images. We refer to this as the task and scene gap (G1).

The second issue concerns hyperparameters. The accuracy of a detector depends on training settings such as the learning rate, momentum, weight decay and the strength of augmentation. The authors of TrashNet, for instance, attributed the 22\% test accuracy of their convolutional neural network (CNN) to poorly chosen hyperparameters~\cite{yang2016classification}. Metaheuristics have been used to tune CNNs, e.g., particle swarm optimization (PSO)~\cite{kennedy1995particle} for a mammography classifier~\cite{aguerchi2024cnn} and a beluga whale optimizer for an InceptionV3 waste classifier~\cite{sayed2024intelligent}. The evidence behind the reported improvements is nevertheless limited. In the studies that we reviewed, metaheuristic tuning in waste vision was applied only to classifiers, and we did not find any waste detector whose training hyperparameters were tuned by a swarm optimizer (G2). Furthermore, the tuned models were compared either with other metaheuristics or with no baseline at all, but never with random search~\cite{bergstra2012random} at an equal number of evaluations, and the test results were given as single values without confidence intervals (G3). Under these conditions it is not possible to say whether an improvement is due to the optimizer or simply to the fact that many settings were tried.

The data form the third issue. According to recent reviews, many waste-sorting studies work in simplified environments with artificially collected images~\cite{lu2022computer}, which do not resemble real waste~\cite{fang2023artificial}. Accuracy also varies strongly from one dataset to another. The same YOLOv5x model, trained and tested separately on two datasets, obtained 84.9 AP$_{50}$ on PlastOPol but only 63.3 on TACO~\cite{cordova2022litter}. An improvement obtained by tuning on studio images may therefore not carry over to street photographs. This matters for India, where segregation at source is compulsory. The closest available resource, the Dense Waste Segmentation Dataset (DWSD)~\cite{ali2025dwsd}, reports only pixel accuracy for segmentation, and we could not find an Indian real-world waste test set with published object-detection baselines (G4).

In this work, named SwarmSort, the tuning question is studied under controlled conditions. We train YOLOv8n~\cite{jocher2023ultralytics}, a nano-sized detector that can run on a smart bin or a mobile phone, and compare three ways of choosing six of its training hyperparameters, which we call arms: the library defaults (Arm~A), the best of 32 random samples (Arm~B) and our own implementation of PSO with 8 particles and 4 rounds (Arm~C). Each search is allowed 32 short proxy trainings. The setting chosen by each arm is trained once for 100 epochs and tested once on a held-out split, and every difference is reported with a paired-bootstrap 95\% confidence interval. The experiment was conducted in two generations. \genone used studio-style images of six classes, and \gentwo added real-world photographs, a seventh class (OTHER) and an Indian test set derived from DWSD. The following research questions are addressed:
\begin{enumerate}
\item[RQ1:] With an equal budget of 32 proxy trainings, does PSO find YOLOv8n training hyperparameters that perform better than those of random search and the default settings on a held-out test set? (G2, G3)
\item[RQ2:] Does a short proxy training (12 epochs on 40\% of the data) rank hyperparameter settings in the same order as a longer training? (G3)
\item[RQ3:] Does the answer to RQ1 change when the data move from studio images to real-world and Indian scenes, and how large is the drop from studio to real-world images? (G1, G4)
\end{enumerate}

The main contributions of this paper are summarized below.
\begin{itemize}
\item \textit{A lightweight multi-object waste detector} (G1). YOLOv8n is trained on multi-object studio and real-world scenes, with six classes in \genone and seven in \gentwo (Sections~\ref{subsec:datasets} and~\ref{subsec:detector}).
\item \textit{Swarm-based tuning of a waste detector} (G2). A PSO implemented by us searches the learning-rate schedule, momentum, weight decay and two augmentation settings (Sections~\ref{subsec:space} and~\ref{subsec:pso}). To the best of our knowledge, the training hyperparameters of a waste detector have not been tuned with a swarm optimizer before.
\item \textit{An equal-budget protocol with three arms} (G3). The protocol includes a proxy-validity check whose pass threshold was decided beforehand (passed with Spearman $\rho = 0.976$, $n = 8$, \genone only), a test split that is used only once, and paired-bootstrap 95\% confidence intervals on all differences (Sections~\ref{subsec:design}, \ref{subsec:proxy} and~\ref{subsec:eval}).
\item \textit{Evaluation over two generations and three domains} (G4). Studio, real-world and Indian results are reported separately. These include detection results on DWSD, for which we are not aware of any published detection baseline. It should be noted that our DWSD boxes are derived from mask regions in which touching items are merged, and that the class mapping had to be inferred from the masks, so these numbers also reflect a mismatch in labeling. On the same test images, real-world training data increased real-world \mapfifty from 0.038--0.096 (\genone) to 0.338--0.353 (\gentwo), with a loss of 0.074--0.105 on studio images; all these changes are significant (Section~\ref{subsec:res-domain}).
\item \textit{Reporting of null results} (O5). PSO did not beat random search by any meaningful margin. Its only significant advantage, $+0.003$ \mapfifty on the Indian set, occurs where all scores are close to zero. Tuning gave no gain over the defaults in \genone and a small gain in \gentwo ($+0.017$ and $+0.020$ on the mixed test set), which did not extend to real-world or Indian photographs. We discuss the reasons in Section~\ref{sec:discussion}, and we release the code, search logs, result files and trained models.
\end{itemize}

The rest of the paper is organized as follows. Section~\ref{sec:related} reviews the related work and identifies the gaps G1--G4. Section~\ref{sec:methods} describes the datasets and methods. Section~\ref{sec:results} presents the results for RQ2, RQ1 and RQ3, in that order. Section~\ref{sec:discussion} discusses the findings along with their implications and limitations, Section~\ref{sec:future} outlines future work, and Section~\ref{sec:conclusion} concludes the paper.
